\documentclass{article}
% ---------------------------------------------------------------
% Póster de revisión de literatura (A0 vertical)
% Finanzas I (055D), UNCP, 2026-II. Compilar con pdfLaTeX.
% ---------------------------------------------------------------
\usepackage[paperwidth=841mm,paperheight=1189mm,margin=0mm]{geometry}
\usepackage[utf8]{inputenc}
\usepackage[T1]{fontenc}
\IfFileExists{spanish.ldf}{\usepackage[spanish,es-noshorthands,es-nodecimaldot]{babel}}{}
\usepackage{helvet}
\renewcommand{\familydefault}{\sfdefault}
\usepackage[letterspace=120]{microtype}
\usepackage{anyfontsize}
\usepackage{graphicx}
\usepackage{xcolor}
\usepackage[most]{tcolorbox}
\usepackage{enumitem}
\usepackage{multicol}
\usepackage{xurl}
\usepackage[hidelinks]{hyperref}
\usepackage{qrcode}
\urlstyle{same}
\pagestyle{empty}
\setlength{\parindent}{0pt}
\hyphenpenalty=10000 \exhyphenpenalty=10000

% ---------------- Paleta ----------------
\definecolor{navy}{HTML}{8E0C11}
\definecolor{navyLight}{HTML}{B8121A}
\definecolor{gold}{HTML}{00A3E0}
\definecolor{acento}{HTML}{A6E1F7}
\definecolor{goldSoft}{HTML}{E8F6FC}
\definecolor{teal}{HTML}{353435}
\definecolor{tealSoft}{HTML}{F1EEEE}
\definecolor{ink}{HTML}{231F20}
\definecolor{muted}{HTML}{5E5757}
\definecolor{line}{HTML}{E5D7D7}
\definecolor{page}{HTML}{F6F1F1}
\definecolor{mist}{HTML}{F7D3D5}
\pagecolor{page}

% ---------------- Tamaños (ajustados por póster) ----------------
\newcommand{\bodysize}{\fontsize{30}{39.6}\selectfont}
\newcommand{\cardtitlesize}{\fontsize{35.4}{42.5}\selectfont}
\newcommand{\smallsize}{\fontsize{24.6}{32.0}\selectfont}

% ---------------- Estilos de tarjeta ----------------
\tcbset{
  card/.style={enhanced, sharp corners, colback=white, colframe=line, boxrule=0.9mm,
    left=14mm, right=14mm, top=12mm, bottom=12mm, fontupper=\bodysize\color{ink}\raggedright,
    overlay={\fill[#1] (frame.north west) rectangle ([yshift=-5mm]frame.north east);}},
  card/.default=navy,
}
\newcommand{\cardtitle}[2][navy]{{\color{#1}\cardtitlesize\bfseries\textls[110]{\MakeUppercase{#2}}}\par\vspace{7mm}}
\newcommand{\bul}{\raisebox{0.18em}{\color{gold}\rule{0.5em}{0.5em}}}
\setlist[itemize]{label=\bul, leftmargin=1.25em, labelsep=0.6em, itemsep=0.55em, topsep=0pt, parsep=0pt}
\newcommand{\pill}[3]{\tcbox[on line, enhanced, arc=7mm, boxrule=0.8mm, colback=#1, colframe=#2,
  left=7mm, right=7mm, top=3mm, bottom=3mm, fontupper=\fontsize{26}{30}\selectfont\bfseries]{#3}}

\begin{document}

% ======================= ENCABEZADO =======================
\begin{tcolorbox}[enhanced, sharp corners, colback=navy, boxrule=0pt, width=\paperwidth,
  left=34mm, right=34mm, top=20mm, bottom=20mm,
  overlay={\fill[gold] (frame.south west) rectangle ([yshift=5mm]frame.south east);
           \fill[navyLight] ([xshift=-160mm]frame.north east) -- (frame.north east) -- ([yshift=-160mm]frame.north east) -- cycle;}]
\begin{minipage}[c]{92mm}\centering\includegraphics[width=\linewidth]{Escudo_UNCP.png}\end{minipage}\hfill
\begin{minipage}[c]{555mm}\centering
{\color{acento}\fontsize{25}{30}\selectfont\bfseries\textls[160]{REVISIÓN DE LITERATURA \quad|\quad FINANZAS I \quad|\quad TEMA 61}\par}
\vspace{10mm}
{\color{white}\fontsize{52}{58.8}\selectfont\bfseries Testing the pecking order theory of capital structure: Evidence from Turkey using panel quantile regression approach\par}
\vspace{8mm}
{\color{mist}\fontsize{34}{41.5}\selectfont\itshape Prueba de la teoría de la jerarquía financiera de la estructura de capital: evidencia de Turquía con regresión por cuantiles en panel\par}
\vspace{10mm}
{\color{white}\fontsize{29}{35}\selectfont\bfseries Durmuş Yıldırım y Ali Kemal Çelik\par}
\vspace{9mm}
\pill{white}{white}{\color{navy}Q1}\hspace{6mm}\pill{navy}{white}{\color{white}\mdseries \textit{Borsa Istanbul Review}, 21(4), 317-331 (2021)}
\end{minipage}\hfill
\begin{minipage}[c]{92mm}\centering\includegraphics[width=\linewidth]{Logo_FE.png}\end{minipage}
\end{tcolorbox}
\vspace{-1mm}
\begin{tcolorbox}[enhanced, sharp corners, colback=goldSoft, boxrule=0pt, width=\paperwidth,
  left=34mm, right=34mm, top=8mm, bottom=8mm, halign=center]
{\color{navy}\fontsize{27}{36}\selectfont \textbf{Póster elaborado por:} Meglinho Sneyder Rojas Poma \enspace\textbar\enspace Código 2024200522B\par \fontsize{24}{32}\selectfont Escuela Profesional de Economía, Universidad Nacional del Centro del Perú \enspace\textbar\enspace Finanzas I (055D), 2026-II\par}
\end{tcolorbox}

\vfill
% ======================= CUERPO =======================
\hspace*{34mm}\begin{minipage}{773mm}
\begin{tcolorbox}[enhanced, sharp corners, colback=navy, boxrule=0pt, left=16mm, right=16mm, top=12mm, bottom=11mm,
  fontupper=\bodysize\raggedright, overlay={\fill[gold] (frame.north west) rectangle ([yshift=-5mm]frame.north east);}]
\begin{minipage}[t]{237.7mm}\raggedright{\color{acento}\smallsize\bfseries\textls[110]{\MakeUppercase{Mercado}}}\par\vspace{2mm}{\color{white}Manufactureras de Borsa Istanbul}\end{minipage}\hfill\begin{minipage}[t]{237.7mm}\raggedright{\color{acento}\smallsize\bfseries\textls[110]{\MakeUppercase{Periodo}}}\par\vspace{2mm}{\color{white}2000 a 2018}\end{minipage}\hfill\begin{minipage}[t]{237.7mm}\raggedright{\color{acento}\smallsize\bfseries\textls[110]{\MakeUppercase{Método}}}\par\vspace{2mm}{\color{white}Regresión por cuantiles en panel}\end{minipage}
\end{tcolorbox}
\end{minipage}

\vfill
\hspace*{34mm}\begin{minipage}{773mm}
\begin{tcbraster}[raster columns=3, raster equal height, raster column skip=16mm]
\begin{tcolorbox}[card]\cardtitle{Contexto}La teoría de la jerarquía financiera dice que las empresas prefieren financiarse primero con fondos propios, luego con deuda y al final con acciones, porque emitir acciones es lo más caro.\end{tcolorbox}
\begin{tcolorbox}[card]\cardtitle{Problema}No se sabía si esa jerarquía se cumple igual en distintos niveles de inversión de las empresas manufactureras turcas.\end{tcolorbox}
\begin{tcolorbox}[card]\cardtitle{Objetivo}Probar la teoría de la jerarquía financiera en distintos niveles de inversión para las manufactureras de Borsa Istanbul.\end{tcolorbox}
\end{tcbraster}
\end{minipage}

\vfill
\hspace*{34mm}\begin{minipage}{773mm}
\begin{tcolorbox}[card]\cardtitle{Metodología}
\begin{minipage}[t]{0.48\linewidth}
\begin{itemize}
\item Se mide la importancia de los fondos internos y externos para financiar la inversión.
\item Se estima en distintos cuantiles del nivel de inversión.
\item Se compara por tamaño, apalancamiento, periodo y subsector.
\end{itemize}
\end{minipage}\hfill
\begin{minipage}[t]{0.48\linewidth}
\begin{itemize}
\item En el panel A solo se controla por la razón valor de mercado sobre valor contable.
\item En el panel B se añaden tamaño (logaritmo de ventas), caja y apalancamiento.
\end{itemize}
\end{minipage}
\end{tcolorbox}
\end{minipage}

\vfill
\hspace*{34mm}\begin{minipage}{773mm}
\begin{tcolorbox}[card]\cardtitle{Figura clave del estudio}
\begin{minipage}[c]{420.0mm}\centering
\fcolorbox{line}{white}{\includegraphics[width=\dimexpr420.0mm-2\fboxsep-2\fboxrule\relax]{figura.png}}
\end{minipage}\hfill
\begin{minipage}[c]{\dimexpr\linewidth-420.0mm-18mm\relax}\raggedright
{\color{navy}\bfseries Figura 2.}\enspace{\itshape Coeficientes de la regresión cuantílica del primer escalón de la jerarquía financiera: sensibilidad de la inversión a los fondos internos (CF/K) y externos (EXTF/K) por decil de inversión. Panel A: control por M/B; panel B: con controles adicionales.}\par\vspace{8mm}
{\color{teal}\smallsize\bfseries\textls[110]{CÓMO LEERLA}}\par\vspace{4mm}
\begin{itemize}
\item La línea punteada son los fondos internos y la discontinua los fondos externos.
\item En el panel A ambas pendientes crecen con el nivel de inversión, pero más la de fondos externos: de 0.035 a 0.081 entre el primer y el noveno decil.
\item La sensibilidad a los fondos internos sube solo un poco, de 0.084 a 0.100.
\end{itemize}
\vspace{8mm}
{\color{muted}\smallsize Figura tomada de Yıldırım y Çelik (2021).}
\end{minipage}
\end{tcolorbox}
\end{minipage}

\vfill
\hspace*{34mm}\begin{minipage}{773mm}
\begin{tcolorbox}[card=teal]\cardtitle[teal]{Resultados}
\begin{tcolorbox}[enhanced, sharp corners, colback=tealSoft, boxrule=0pt, left=14mm, right=14mm, top=10mm, bottom=11mm,
  borderline west={4mm}{0mm}{teal}, fontupper=\bodysize\color{ink}\raggedright]
{\color{teal}\smallsize\bfseries\textls[110]{HALLAZGO CENTRAL}}\par\vspace{4mm}
{\bfseries La teoría de la jerarquía financiera se cumple, y la sensibilidad a fondos internos y a la deuda crece con el nivel de inversión.}
\end{tcolorbox}
\vspace{11mm}
\begin{minipage}[t]{0.48\linewidth}
\begin{itemize}
\item Las empresas pequeñas siguen la jerarquía.
\item No se cumple en empresas muy o poco apalancadas.
\item Antes de la crisis (2000 a 2009) se siguió la jerarquía; después (2010 a 2018) se prefirió la deuda.
\end{itemize}
\end{minipage}\hfill
\begin{minipage}[t]{0.48\linewidth}
\begin{itemize}
\item Alimentos y bebidas, químicos y minerales no metálicos siguen la jerarquía.
\item Las empresas usan primero los fondos internos y luego los externos, como predice la jerarquía financiera.
\item Al aumentar el nivel de inversión el orden de preferencia no cambia.
\end{itemize}
\end{minipage}
\end{tcolorbox}
\end{minipage}

\vfill
\hspace*{34mm}\begin{minipage}{773mm}
\begin{tcbraster}[raster columns=3, raster equal height, raster column skip=16mm]
\begin{tcolorbox}[card, raster multicolumn=2]\cardtitle{Conclusiones}Las empresas usan primero sus fondos internos y luego recurren a fondos externos, aunque la preferencia cambia con el apalancamiento y con el periodo económico.\end{tcolorbox}
\begin{tcolorbox}[card=gold, colback=goldSoft]\cardtitle{Aporte al tema 61}Muestra por qué el costo del patrimonio importa: si emitir acciones es caro, las empresas prefieren otras fuentes. El tema 61 estima ese costo para consumo y construcción, sectores parecidos a los del estudio.\end{tcolorbox}
\end{tcbraster}
\end{minipage}

\vfill
% ======================= PIE =======================
\begin{tcolorbox}[enhanced, sharp corners, colback=navy, boxrule=0pt, width=\paperwidth,
  left=34mm, right=34mm, top=14mm, bottom=11mm,
  overlay={\fill[gold] (frame.north west) rectangle ([yshift=-5mm]frame.north east);}]
\begin{minipage}[c]{245mm}\raggedright
{\color{acento}\smallsize\bfseries\textls[110]{FUENTE DE DATOS DEL ESTUDIO}}\par\vspace{5mm}
{\color{white}\smallsize Estados financieros de empresas manufactureras de Borsa Istanbul, de 2000 a 2018.}
\end{minipage}\hfill
\begin{minipage}[c]{385mm}\raggedright
{\color{acento}\smallsize\bfseries\textls[110]{REFERENCIA}}\par\vspace{5mm}
{\color{white}\smallsize\hangindent=1.2em\hangafter=1 Yıldırım, D., y Çelik, A. K. (2021). Testing the pecking order theory of capital structure: Evidence from Turkey using panel quantile regression approach. \textit{Borsa Istanbul Review}, \textit{21}(4), 317-331. \url{https://doi.org/10.1016/j.bir.2020.11.002}\par}
\end{minipage}\hfill
\begin{minipage}[c]{100mm}\centering
\tcbox[enhanced, sharp corners, colback=white, boxrule=0pt, left=3mm, right=3mm, top=3mm, bottom=3mm]{\qrcode[hyperlink, height=46mm]{https://doi.org/10.1016/j.bir.2020.11.002}}\par\vspace{4mm}
{\color{acento}\fontsize{17}{21}\selectfont\bfseries\textls[100]{ESCANEA Y LEE EL ARTÍCULO}}
\end{minipage}
\end{tcolorbox}
\end{document}
